\documentclass[11pt,a4paper]{article}
\usepackage[utf8]{inputenc}
\usepackage[T1]{fontenc}
\usepackage{lmodern}
\usepackage{textcomp}
\usepackage{amsmath,amssymb}
\usepackage{graphicx}
\usepackage{booktabs}
\usepackage{longtable}
\usepackage{array}
\usepackage{listings}
\usepackage{xcolor}
\usepackage[margin=1in]{geometry}
\usepackage{float}
\usepackage[hidelinks,breaklinks]{hyperref}
\usepackage{microtype}
\setlength{\parskip}{4pt}
\lstset{basicstyle=\ttfamily\scriptsize,breaklines=true,breakatwhitespace=false,
  postbreak=\mbox{\textcolor{gray}{$\hookrightarrow$}\space},frame=single,
  framerule=0.3pt,rulecolor=\color{gray!60},columns=fullflexible,keepspaces=true,
  showstringspaces=false,numbers=none,xleftmargin=4pt,xrightmargin=4pt,
  literate={\_}{{\_}}1}
\graphicspath{{figures/}}
\providecommand{\tightlist}{}
\title{The Hubble Tension: A Conceptual Summary of Early- and Late-Universe Determinations of $H_0$}
\author{HEP-Genesis-Agent session record \\ \small Project: ramp-features}
\date{2026-10-15}
\begin{document}
\maketitle
\noindent\textit{\small This document is an AI-generated account of an interactive session with the HEP-Genesis agent. Its text was written by a language model strictly from the session record; figures and numbers come only from tool results obtained during the session. The supplementary material and the transcript, generated alongside, document exactly what ran. Verify independently before use.}
\vspace{1em}

\begin{abstract}
This report records a short, conceptual session in which the user asked what the Hubble tension is, requesting a two-paragraph explanation followed by a written report. The assistant described the tension as a disagreement between two kinds of measurement of the present-day expansion rate $H_0$. Fitting $\Lambda$CDM to the cosmic microwave background gives $H_0 = 67.4 \pm 0.5$ km\,s$^{-1}$\,Mpc$^{-1}$ (Planck, arXiv:1807.06209). The Cepheid-calibrated Type Ia supernova distance ladder gives $H_0 = 73.04 \pm 1.04$ km\,s$^{-1}$\,Mpc$^{-1}$ (SH0ES, arXiv:2112.04510). The assistant stated that the two differ at about $5\sigma$. It also summarized two classes of explanation, measurement systematics and physics beyond $\Lambda$CDM. In this context it cited an intermediate tip-of-the-red-giant-branch calibration near $69.8$ km\,s$^{-1}$\,Mpc$^{-1}$ (arXiv:1907.05922) and a review of proposed models (arXiv:2103.01183). No computation, data analysis, or literature search was run in the session. All values and citations come from the assistant's own knowledge, as the assistant itself noted, and none was checked against the sources during the session.
\end{abstract}

\section{Introduction}
The session contained one question: ``What is the Hubble tension?'' The user asked for two paragraphs of explanation, followed by a report. The scope is therefore conceptual. The aim was to state what quantity is in dispute, which measurements disagree and by how much, and what explanations have been proposed. No new measurement, fit, or reanalysis was requested or carried out. This report sets out the content of the assistant's explanation and makes clear how much evidence supports it.

\section{Methods}
No tools were used. The session ran in a quick-response mode. No MCP servers, HPC jobs, code, or literature-search tools were called, and no figures were produced. The content below is the assistant's written explanation, drawn from its own knowledge. The assistant said explicitly that its citations ``come from my own knowledge, not from a literature search in this session'' and that the values are ``as published in the papers named.'' Because nothing was retrieved, none of the numbers or citations was checked against the original papers during the session. The accompanying supplement contains the session record. There are no tool calls or code to reproduce.

\section{Results}

\subsection{Definition of the tension}
The assistant defined the Hubble tension as a disagreement between two families of methods for measuring the current expansion rate, the Hubble constant $H_0$.

The first family consists of early-universe, model-dependent determinations. The standard $\Lambda$CDM model, with dark energy and cold dark matter, is fitted to the cosmic microwave background (CMB) and then evolved forward to the present epoch. By the assistant's account, the Planck satellite obtains
\begin{equation}
H_0^{\mathrm{CMB}} = 67.4 \pm 0.5\ \mathrm{km\,s^{-1}\,Mpc^{-1}}
\end{equation}
in this way (Planck Collaboration 2018, arXiv:1807.06209). The assistant added that baryon acoustic oscillations (BAO) and Big Bang nucleosynthesis give a similar value. It did not quote specific numbers for these.

The second family consists of late-universe distance-ladder determinations. The ladder has three rungs:
\begin{enumerate}
\item parallaxes and other geometric anchors calibrate Cepheid variable stars;
\item the Cepheids calibrate Type Ia supernovae;
\item the supernovae measure distances deep into the Hubble flow.
\end{enumerate}
The SH0ES result quoted by the assistant is
\begin{equation}
H_0^{\mathrm{ladder}} = 73.04 \pm 1.04\ \mathrm{km\,s^{-1}\,Mpc^{-1}}
\end{equation}
(Riess et al. 2022, arXiv:2112.04510).

The assistant stated that these values differ by about $5\sigma$. It argued that a gap this large is too big to attribute to chance, provided the quoted uncertainties are correct. This significance figure is the assistant's statement and was not recomputed in the session.

\subsection{Classes of proposed explanation}
The assistant described two broad classes of explanation and said that neither has prevailed.

\paragraph{Measurement systematics.} The assistant listed several candidate systematic errors:
\begin{itemize}
\item crowding and photometry problems in Cepheid measurements;
\item the dependence of Cepheid luminosity on metallicity;
\item the procedure used to standardize Type Ia supernovae.
\end{itemize}
As evidence that this debate remains open, it cited a distance ladder calibrated with the tip of the red giant branch (TRGB) instead of Cepheids. That calibration gives an intermediate value near $69.8$ km\,s$^{-1}$\,Mpc$^{-1}$ (Freedman et al. 2019, arXiv:1907.05922).

\paragraph{Physics beyond $\Lambda$CDM.} The assistant listed the following proposals:
\begin{itemize}
\item early dark energy;
\item additional relativistic species;
\item modifications to recombination;
\item interactions in the dark sector.
\end{itemize}
It characterized most of these as easing the $H_0$ discrepancy only at the cost of worse fits to other data. The examples it gave were galaxy clustering, as quantified by $S_8$, and BAO. For this point it referred to the review by Di Valentino et al. 2021 (arXiv:2103.01183).

\subsection{Prospective independent probes}
The assistant identified three independent approaches that may settle the issue:
\begin{itemize}
\item gravitational-wave ``standard sirens'';
\item time delays in strongly lensed quasars;
\item new JWST observations of Cepheids and red giants.
\end{itemize}
It gave no numerical results from these methods.

\section{Discussion}
The session's explanation follows a clear logic. The tension is a disagreement between an inference that depends on a cosmological model, extrapolated from the CMB epoch, and a measurement built from successive empirical calibrations in the local universe. The location of the TRGB-based value between the Planck and SH0ES results matters in that framing. It suggests that the choice of calibrator in the ladder affects the result, which is why the assistant treated systematics as a live explanation. The assistant's description of new-physics models presents the counterpart difficulty. A modification that raises the CMB-inferred $H_0$ tends to make the fit worse for other observables, such as $S_8$ and BAO, so these proposals do not resolve the tension without cost.

The main limitation is evidential. Every number and citation in this report was supplied from the assistant's knowledge, not retrieved or checked during the session, and the assistant said so. The $5\sigma$ significance was asserted and not computed. The assistant did not specify how the uncertainties were combined or whether any correlations were assumed. The BAO and nucleosynthesis agreement with Planck was stated qualitatively, without values. The descriptions of the systematic effects and new-physics models are summaries and were not tested quantitatively. Readers should therefore treat the content as a conceptual overview whose figures need to be confirmed against the cited papers.

\section{Conclusions}
The session answered the user's conceptual question. According to the assistant, the Hubble tension is a discrepancy of about $5\sigma$ between two determinations of $H_0$. The CMB-based $\Lambda$CDM inference from Planck gives $67.4 \pm 0.5$ km\,s$^{-1}$\,Mpc$^{-1}$. The SH0ES Cepheid--supernova distance ladder gives $73.04 \pm 1.04$ km\,s$^{-1}$\,Mpc$^{-1}$. A TRGB calibration gives an intermediate value near $69.8$ km\,s$^{-1}$\,Mpc$^{-1}$.

Two explanations remain open: measurement systematics, and physics beyond $\Lambda$CDM. Proposals of the second kind generally worsen agreement with $S_8$ or BAO. The assistant named standard sirens, lensed-quasar time delays, and JWST Cepheid and red-giant observations as possible ways to decide between the explanations.

No results were computed or checked against sources in the session. A natural next step, implied by the assistant's caveat, would be to retrieve the cited papers and confirm the quoted values.

\section*{References}
\begin{itemize}
\item Planck Collaboration 2018, arXiv:1807.06209.
\item Riess et al. 2022, arXiv:2112.04510.
\item Freedman et al. 2019, arXiv:1907.05922.
\item Di Valentino et al. 2021 (review), arXiv:2103.01183.
\end{itemize}

\section*{Data and code availability}
All outputs of this session are collected in the project folder \href{file:///Users/nramachandra/hep-genesis-testing/projects/ramp-features}{\texttt{/\allowbreak{}Users/\allowbreak{}nramachandra/\allowbreak{}hep-genesis-testing/\allowbreak{}projects/\allowbreak{}ramp-features}}.
Full documentation of how these results were obtained is in the supplementary material (\texttt{supplement.pdf}: implementation details, the tool-call log, the research board with every finding and its evidence, figures not shown here, and the code index) and the verbatim transcript (\texttt{transcript.pdf}), generated alongside this report.


\end{document}
